Our findings highlight the superiority of encoder-based models over decoder-based LLMs for event forecasting, emphasizing the limitations of autoregressive models in structured prediction tasks. Among encoder models, \textit{BioClinical-ModernBERT-base} consistently performs best, achieving superior concordance (0.677 vs 0.646) and better external generalization despite slightly lower internal F1 scores than general-purpose \textit{ModernBERT-base}. MLP-head fine-tuning excels in F1 scores while fine-tuned masking models achieve higher concordance, particularly in long-term forecasting. This suggests that classification and ranking tasks benefit from distinct optimization strategies.

The consistent advantage of biomedically-pretrained models across tasks deserves special attention. While general-purpose models achieve strong internal performance, biomedical pretraining provides crucial benefits for real-world deployment: \textit{BioClinical-ModernBERT} variants show the smallest performance degradation on external datasets and maintain higher zero-shot capabilities (0.246 F1 at 1h vs near-zero for standard BERT/RoBERTa). This suggests that domain-specific pretraining enhances not just accuracy but also robustness and generalization—critical factors when considering clinical deployment where distribution shifts are common.

The gradual improvement of F1 scores over longer time horizons indicates that event patterns become more predictable over time, whereas short-term forecasting remains challenging due to higher variability. Performance drops on external validation datasets highlight generalization challenges, though \textit{ModernBERT} models exhibit relative robustness. For survival analysis, instruction-tuned LLMs like \textit{Llama-3.3-70B-Instruct} outperform traditional transformer baselines, with several models achieving peak concordance at early timepoints rather than with extended observation.

Our sensitivity analysis reveals several trade-offs: time-ordered training generally improves concordance while text-ordering can yield better F1 scores on externally-annotated datasets. Robustness experiments show F1 scores degrade significantly beyond 60\% timestep dropout, while concordance remains stable, indicating event ranking is less sensitive to partial history than event classification.

\textbf{Methodological Contributions and Potential Impact.} Our framework demonstrates how narrative clinical texts can be systematically converted into structured temporal representations for forecasting tasks. The ability to extract temporally structured insights from unstructured clinical text could potentially support clinical decision-making, particularly in settings with limited access to structured data or specialized expertise. 

Importantly, while consumers frequently consult LLMs about health risk via prompting in chats, this study demonstrates that the prompting approach performs substantially worse in risk prediction than alternative approaches, at least with respect to precision/recall/F1, with prompted LLMs achieving at best 0.460 F1 at 24h compared to 0.653 for fine-tuned encoders.
Our study highlights several alternative approaches and characterizes their relative performance strengths and weaknesses.
In addition, our finding that concordance remains stable with 60\% missing context suggests robustness to incomplete documentation scenarios, which highlights the degradation pattern in performance as it relates to context availability. 

Our approach also demonstrates that our language model systems can reliably capture the temporal reasoning that clinicians use in practice based on their clinical documentation. However, additional work would be required to validate performance on real-time clinical data and demonstrate measurable impact on patient outcomes in deployment. Further discussion on societal impact of our work is in Appendix \ref{apd:broader_impact}.

\textbf{Limitations and Future Directions.} 
Our study has important limitations to consider. First, the pipeline relies on case reports from the PubMed Open Access (PMOA) corpus, which often reflect rare or atypical presentations and differ from routine clinical notes (e.g., progress notes, discharge summaries), potentially limiting the model's generalizability to real-world health system texts. Compared to routine clinical notes, case reports offer rich, temporally structured narratives with explicit clinical rationale, making them especially useful for evaluating model performance under sparse documentation. This study develops a general framework for textual time series forecasting, using case reports as an interpretable, temporally rich foundation for method development with potential extensions to real-world clinical corpora like MIMIC discharge summaries.

Second, while we focused on sepsis due to its clinical relevance and prevalence in case reports, our framework is fundamentally disease-agnostic. Preliminary results from over 125K PMOA case reports across various diagnostic conditions show promising generalization, with future work planned on broader diagnostic categories and real-world corpora such as MIMIC-IV discharge summaries.

Third, despite efforts to mitigate causal leakage during training and evaluation (e.g., timestamp extraction, temporal masking, and use of external test sets), pretrained language models may have been exposed to PMOA content. Although these models were not trained on our derived (event, time) sequences or temporally reordered data, it is possible that the original case report narratives were included in their pretraining data. Our pipeline applies a series of transformations—including temporal reordering and masking of future events beyond a prediction cutoff—which fundamentally alter the input representation and enforce a causal structure that is not present in the original documents. This creates a distinct prediction setup that requires models to reason over partial, temporally grounded inputs rather than complete narratives. To further assess the impact of pretraining exposure, we include results from \textit{RedPajama-INCITE-7B-Instruct}, an open-source model with a fully documented, PubMed-free training corpus. Its comparable performance to other decoder models (0.352 F1 at 24h) and similar concordance (0.618) to models potentially exposed to PubMed provides evidence that our T2S2 prediction task is sufficiently distinct from raw PubMed text, helping to mitigate causal leakage concerns. A detailed discussion on the scope of using open-source models for avoiding potential data leakage has been presented in Appendix \ref{apd:open_source}.

Finally, our analysis is limited by the modest size of the high-cost manual annotations collected as a gold standard to calibrate and validate our LLM extraction pipeline. Expanding gold-standard validation across diverse diagnoses could assess the generalizability of our approach and aid in evaluating model robustness, and is a direction for future work.